\documentclass[letterpaper]{article}

\usepackage[letterpaper,margin=1in]{geometry}

% Compile from the repository root so this path resolves to the pinned WSP
% submodule.
\input{wsp/testing/test-case-library.tex}

\begin{document}
\def\TCID{TC-0008}
\def\TCTitle{SOUND and PLAY}
\def\TCPurpose{Verify the tones SOUND and PLAY make: the range of SOUND, the
notes, lengths, dots, pauses, octaves, tempo and styles of the music macro
language, the state that lasts from one PLAY to the next, and error 5.}
\def\TCPriority{\TCPriorityHigh}

\def\TCRequirementRef{REQ-0007}
\def\TCDesignRef{Not applicable.}
\def\TCFeatureRef{\texttt{include/wfc/brun71/sound.hpp};
\texttt{src/sound.cpp}; \texttt{tests/sound\_test.cpp}.}
\def\TCTestTechnique{Equivalence classes of the commands; boundary values of
the frequency, duration, octave, length, tempo and note numbers.}

\def\TCPreconditions{The library is built and the test program
\texttt{brun71\_sound\_test} exists.}
\def\TCEnvironment{Any supported target; no sound is played. CTest runs the
program with a 60 second timeout.}
\def\TCAssumptions{The note frequencies (N46 is 440 Hz) and the styles
(7/8, all, 3/4 of the length) of REQ-0007, not measured on BRUN71.}

\def\TCInputData{SOUND arguments at and beyond their limits; music strings
with each command, upper and lower case, blanks; refused strings.}
\def\TCInitialState{A new player: O4, L4, T120, MN.}

\def\TCProcedure{\begin{enumerate}
  \item Run \texttt{brun71\_sound\_test} (CTest \texttt{TC-0008-sound}, in
    WFC \texttt{wfc-brun71-TC-0008-sound}).
  \item Compare the tones (frequency and milliseconds) of each SOUND and
    music string with the expected ones.
  \item Check the octave, length and tempo kept after each string.
  \item Check that the refused arguments and strings raise error 5.
\end{enumerate}}
\def\TCExpectedResult{The program exits with status 0 and prints
\texttt{sound: ok}; no \texttt{FAIL} line is written.}
\def\TCPostConditions{None; the test writes no files.}

\TCTemplate

\end{document}
